\documentclass[11pt]{article}

\usepackage{amsmath,amssymb,amsthm,mathtools}
\usepackage[T1]{fontenc}
\usepackage{lmodern}
\usepackage[margin=1in]{geometry}
\usepackage{enumitem}
\usepackage{hyperref}
\usepackage{microtype}
\usepackage{xcolor}

\hypersetup{colorlinks=true,linkcolor=blue,citecolor=blue,urlcolor=blue}

\newtheorem{theorem}{Theorem}[section]
\newtheorem{proposition}[theorem]{Proposition}
\newtheorem{lemma}[theorem]{Lemma}
\newtheorem{corollary}[theorem]{Corollary}
\newtheorem{remark}[theorem]{Remark}
\newtheorem{definition}[theorem]{Definition}
\newtheorem{placeholder}[theorem]{Placeholder}

\title{D5 131: what can be proved now about the front-end seed statement \textup{032}}
\author{}
\date{March 2026}

\begin{document}
\maketitle

\begin{abstract}
This note revisits the front-end seed-isolation statement usually called
\textup{032}.
The current accessible bundle still does not contain the original endpoint-word
catalog from the old front-end artifact, so a completely first-principles proof
of seed isolation is not yet available from the present files alone.
What can be proved cleanly now is the following.
Once the live seed class has been reduced to a single \emph{oriented
bridge-type endpoint-word class}---one repeated neighbor on the left, one
shared bridge symbol, and one repeated neighbor on the right---the numerical
representative
\[
L=[2,2,1],\qquad R=[1,4,4]
\]
is forced by the standard front-end normalizations.
So the real remaining burden behind \textup{032} is not a new odd-$m$ argument;
it is the upstream finite endpoint-word classification that the live local seed
class has exactly that bridge form.
\end{abstract}

\section{Scope}

In the revised manuscript, the front-end input \textup{032} is the statement
that, after the earlier tiny static repair branches are removed, the odd-$m$,
$d=5$ endpoint-local front end reduces to the seed
\[
L=[2,2,1],\qquad R=[1,4,4].
\]

The present note does \emph{not} claim a full first-principles derivation of
that statement from the raw endpoint-search artifacts.
Those specific artifact files are not present in the current accessible bundle.
Instead, the note isolates what can already be proved symbolically: the
numerical form $[2,2,1]/[1,4,4]$ is only a \emph{normal form}.
The actual mathematical content that still needs promotion upstream is the much
smaller claim that the live endpoint seed class is a single oriented
bridge-type class.

\section{Abstract seed normal form}

We work only with the direction alphabet
\[
\mathcal D=\{1,2,3,4\},
\]
viewed with its cyclic order
\[
1\to 2\to 3\to 4\to 1.
\]
The accepted front-end symmetries allow cyclic relabeling of the direction
alphabet together with the usual left/right orientation normalization.

\begin{definition}[oriented bridge-type endpoint seed]
An \emph{oriented bridge-type endpoint seed} is a pair of length-$3$ words
\[
L=[\ell,\ell,b],\qquad R=[b,r,r]
\]
in the direction alphabet $\mathcal D$ such that:
\begin{enumerate}[label=\textup{(\roman*)},leftmargin=2.5em]
\item $b$ is the shared bridge symbol between the left and right endpoint
      words;
\item $\ell$ is the distinguished repeated neighbor on the left side of the
      bridge;
\item $r$ is the distinguished repeated neighbor on the right side of the
      bridge;
\item after the standard orientation convention is fixed, $\ell$ is the
      clockwise neighbor of $b$ and $r$ is the counterclockwise neighbor of
      $b$.
\end{enumerate}
\end{definition}

\begin{remark}
This definition is intentionally abstract.
It captures exactly the part of the old front-end seed search that matters for
manuscript packaging: once the live local seed class is known to have one
shared bridge symbol and one repeated oriented neighbor on each side, the
actual digits used to write that class are not additional theorem content.
They are fixed by normalization.
\end{remark}

\begin{proposition}[normal-form lemma]
\label{prop:seed-normal-form}
Every oriented bridge-type endpoint seed is equivalent, under the accepted
front-end relabeling and orientation normalizations, to the unique numerical
representative
\[
L=[2,2,1],\qquad R=[1,4,4].
\]
\end{proposition}

\begin{proof}
Let
\[
L=[\ell,\ell,b],\qquad R=[b,r,r]
\]
be an oriented bridge-type seed.
Apply the accepted cyclic relabeling so that the shared bridge symbol becomes
\[
b=1.
\]
After that relabeling, the left repeated symbol must be the clockwise neighbor
of $1$ in the cyclic order on $\mathcal D$, hence
\[
\ell=2.
\]
Likewise the right repeated symbol must be the counterclockwise neighbor of $1$,
hence
\[
r=4.
\]
Therefore the normalized seed is exactly
\[
L=[2,2,1],\qquad R=[1,4,4].
\]
Uniqueness is immediate because the normalization fixes the bridge symbol and
then the orientation convention fixes which adjacent symbol is written on each
side.
\end{proof}

\begin{corollary}[032 as a normal-form corollary]
\label{cor:032-from-bridge-form}
Assume that, after the earlier tiny static repair branches are pruned, the live
odd-$m$, $d=5$ endpoint-local front end reduces to a single oriented
bridge-type endpoint seed class.
Then the seed statement usually written as \textup{032} follows:
\[
L=[2,2,1],\qquad R=[1,4,4].
\]
\end{corollary}

\begin{proof}
Apply Proposition~\ref{prop:seed-normal-form} to the unique surviving oriented
bridge-type seed class.
\end{proof}

\section{What still remains to close 032 independently}

The previous section proves the \emph{normal-form} part of \textup{032}.
So the true remaining burden is smaller than the old wording
``isolate the best seed.''
It is enough to prove one upstream endpoint-word theorem.

\begin{placeholder}[explicit endpoint-word shape theorem]
\label{ph:shape}
After the earlier tiny static repair branches are removed, prove that the live
odd-$m$, $d=5$ endpoint-local front end reduces to a single oriented
bridge-type endpoint seed class in the sense of Definition~2.1.
Equivalently, prove that the unresolved local seed has exactly one shared bridge
symbol and one repeated oriented neighbor on each side.
\end{placeholder}

\begin{theorem}[what is left of 032]
\label{thm:032-reduced-placeholder}
A fully first-principles proof of \textup{032} is equivalent to
Placeholder~\ref{ph:shape} together with the already-proved normal-form lemma
(Proposition~\ref{prop:seed-normal-form}).
No additional odd-$m$ bookkeeping remains once the endpoint-word shape theorem
is supplied.
\end{theorem}

\begin{proof}
If Placeholder~\ref{ph:shape} holds, then
Corollary~\ref{cor:032-from-bridge-form} gives the normalized seed
$[2,2,1]/[1,4,4]$ immediately.
Conversely, any honest first-principles proof of \textup{032} must determine
what structural shape the live local seed class has before that seed can be
written in normalized coordinates.
That structural classification is exactly Placeholder~\ref{ph:shape}.
\end{proof}

\begin{remark}[practical consequence for the manuscript]
The theorem-side burden behind \textup{032} is therefore smaller than it first
appeared.
The digits $[2,2,1]/[1,4,4]$ are not the hard part.
The hard part is the upstream finite endpoint-word classification that the live
seed class has the oriented bridge form.
So, if later support for the explicit path-level endpoint words is added, there
should be no separate heavy proof step left for \textup{032} itself.
\end{remark}

\section{Status}

From the current accessible bundle one may now say the following precisely.

\begin{enumerate}[label=\textup{(\arabic*)},leftmargin=2.5em]
\item The numerical representative of the best front-end seed is explained:
      it is the canonical normal form of an oriented bridge-type seed class.
\item The full first-principles content of \textup{032} has not yet been
      completely internalized, because the needed endpoint-word shape theorem
      is still external to the current bundle.
\item Therefore \textup{032} should not yet be claimed as a fully proved
      theorem inside the consolidated manuscript, but its remaining gap is now
      sharply localized.
\end{enumerate}

\end{document}
